\chapter{Stakeholders \\
\small{\textit{-- Chan, Moks, Ponte}}
\index{stakeholders} 
\index{Chapter!Stakeholders}
\label{Chapter::Stakeholders}}

This chapter identifies the groups and individuals who will be directly or indirectly impacted by the system. Each stakeholder plays a specific role in ensuring the success, adoption, and continued improvement of the trustworthy natural language interface for quantum optimization.

\section{Primary Stakeholders}

\textbf{End Users} — Individuals who will use the interface to describe optimization problems in natural language and obtain quantum or classical solutions. They include students, researchers, and professionals in fields such as logistics, scheduling, and resource allocation.

\textbf{Developers and Engineers} — Team members and future contributors responsible for implementing and maintaining the translation pipeline, integrating solvers, and ensuring the reliability of the system architecture.

\textbf{Academic Advisors and Sponsors} — Faculty mentors and technical representatives from organizations such as D-Wave who provide research guidance, solver access, and ensure alignment with trustworthy-AI principles.

\section{Secondary Stakeholders}

\textbf{Regulatory and Ethical Review Bodies} — Institutional Review Boards (IRB) or equivalent entities that ensure compliance with data privacy and research ethics, particularly that no personal or sensitive data are used in prompts.

\textbf{Third-Party API Providers} — Organizations such as OpenAI and D-Wave that supply the APIs, SDKs, or solvers upon which the system depends.

\textbf{Competitors and Industry Observers} — Developers of comparable frameworks (e.g., OR-LLM-Agent, LLMOPT) who may use this system’s outcomes as benchmarks for trustworthiness and usability in future research.

\section{Stakeholder Impact Summary}

The system impacts stakeholders by:
\begin{itemize}
    \item Empowering end users to interact with quantum optimization through intuitive natural language input.
    \item Enabling developers to refine modular and verifiable architectures for hybrid quantum-classical systems.
    \item Providing academic sponsors and advisors with measurable research outputs for evaluation and publication.
    \item Supporting ethical compliance and responsible AI development through adherence to data privacy standards.
    \item Encouraging cross-institutional collaboration by releasing prompt libraries and experimental datasets.
\end{itemize}

\bigskip
\noindent
Together, these stakeholders form a collaborative ecosystem that drives innovation, ensures accountability, and extends the practical reach of quantum optimization research.
